\documentclass[10pt]{article}

\usepackage[margin=0.8in]{geometry}
\usepackage{amsmath,amssymb}
\usepackage{xcolor}
\usepackage{enumitem}
\usepackage{booktabs}
\usepackage{microtype}

\definecolor{DiffBlue}{RGB}{0,84,166}
\definecolor{SoftGray}{RGB}{245,247,250}

\newcommand{\Diff}[1]{\textcolor{DiffBlue}{\textbf{#1}}}
\newcommand{\code}[1]{\texttt{#1}}
\newcommand{\Margin}{\Delta}
\newcommand{\Mask}{M}
\newcommand{\Adv}{A}

\setlist[itemize]{leftmargin=1.4em,itemsep=0.2em,topsep=0.25em}
\setlist[enumerate]{leftmargin=1.5em,itemsep=0.2em,topsep=0.25em}

\title{\vspace{-2em}OPD Normalization Strategies: Pseudocode}
\author{}
\date{}

\begin{document}
\maketitle
\vspace{-2em}

\section*{Shared Setup}

For each rollout chunk index $t$, rollout sample $b$, and flattened action token
$k$, OPD first computes the sampled-action teacher margin
\[
  \Margin_{t,b,k}
  =
  \log \pi_{\mathrm{teacher}}(a_{t,b,k}\mid s_{t,b})
  -
  \log \pi_{\mathrm{student,old}}(a_{t,b,k}\mid s_{t,b}).
\]

Let $\Mask_{t,b,k}\in\{0,1\}$ be the valid-token mask. In all modes below, invalid
tokens are zeroed at the end:
\[
  \Adv \leftarrow \Adv \odot \Mask.
\]

For OpenVLA-OFT chunked actions, the last dimension is typically
\[
  k = (h,d), \qquad h\in\{1,\ldots,H\}, \quad d\in\{1,\ldots,D\},
\]
where $H=8$ future actions and $D=7$ action dimensions.

\section*{Mode Summary}

\begin{center}
\begin{tabular}{@{}lll@{}}
\toprule
Mode & What Is Normalized & Dense Token Feedback Preserved? \\
\midrule
\code{group\_zscore} & chunk scalar within rollout group & No \\
\code{token\_zscore} & each flattened token position $k$ & Yes \\
\code{action\_dim\_zscore} & each action dimension $d$ & Yes \\
\code{positive\_clip} & raw positive token margins & Yes \\
\code{teacher\_prob} & sampled teacher probability weights & Yes \\
\bottomrule
\end{tabular}
\end{center}

\section*{1. Current \code{group\_zscore}}

\Diff{Difference: collapse all action tokens in a chunk into one scalar, normalize
that scalar within a rollout group, then broadcast it back to every token.}

\begin{enumerate}
  \item Compute one score per chunk/sample by averaging token margins:
  \[
    s_{t,b}
    =
    \frac{\sum_k \Mask_{t,b,k}\Margin_{t,b,k}}
         {\max\left(1,\sum_k \Mask_{t,b,k}\right)}.
  \]

  \item Partition the batch dimension into groups $g$ of size $G$:
  \[
    b=(g,i), \qquad i\in\{1,\ldots,G\}.
  \]

  \item For every chunk $t$ and group $g$, compute:
  \[
    \mu_{t,g} = \frac{1}{G}\sum_{i=1}^G s_{t,g,i},
    \qquad
    \sigma_{t,g}
    =
    \mathrm{std}_{i=1}^G(s_{t,g,i}).
  \]

  \item Broadcast the normalized scalar to every token in that chunk/sample:
  \[
    \Adv_{t,g,i,k}
    =
    \frac{s_{t,g,i}-\mu_{t,g}}{\sigma_{t,g}+\epsilon}.
  \]
\end{enumerate}

In one line:
\[
  \text{token margins}
  \rightarrow
  \text{chunk mean}
  \rightarrow
  \text{group z-score}
  \rightarrow
  \text{same advantage for all tokens in the chunk}.
\]

\section*{2. \code{token\_zscore}}

\Diff{Difference: keep each flattened token position separate. Do not average the
action chunk into one scalar.}

For each token position $k$:
\[
  \mu_k
  =
  \frac{\sum_{t,b}\Mask_{t,b,k}\Margin_{t,b,k}}
       {\max\left(1,\sum_{t,b}\Mask_{t,b,k}\right)},
\]
\[
  \sigma_k
  =
  \sqrt{
    \frac{\sum_{t,b}\Mask_{t,b,k}(\Margin_{t,b,k}-\mu_k)^2}
         {\max\left(1,\sum_{t,b}\Mask_{t,b,k}\right)}
  }.
\]

Then:
\[
  \Adv_{t,b,k}
  =
  \frac{\Margin_{t,b,k}-\mu_k}{\sigma_k+\epsilon}.
\]

This preserves token-level credit assignment:
\[
  \text{token } k_1 \text{ and token } k_2
  \text{ may receive different OPD feedback inside the same chunk.}
\]

\section*{3. \code{action\_dim\_zscore}}

\Diff{Difference: keep semantic action dimensions separate. Pool over chunk horizon
$h$, rollout chunks $t$, and batch samples $b$, but not over action dimension $d$.}

Reshape:
\[
  \Margin_{t,b,k}
  \rightarrow
  \Margin_{t,b,h,d}.
\]

For each action dimension $d$:
\[
  \mu_d
  =
  \frac{\sum_{t,b,h}\Mask_{t,b,h,d}\Margin_{t,b,h,d}}
       {\max\left(1,\sum_{t,b,h}\Mask_{t,b,h,d}\right)},
\]
\[
  \sigma_d
  =
  \sqrt{
    \frac{\sum_{t,b,h}\Mask_{t,b,h,d}(\Margin_{t,b,h,d}-\mu_d)^2}
         {\max\left(1,\sum_{t,b,h}\Mask_{t,b,h,d}\right)}
  }.
\]

Then:
\[
  \Adv_{t,b,h,d}
  =
  \frac{\Margin_{t,b,h,d}-\mu_d}{\sigma_d+\epsilon}.
\]

This avoids mixing, for example, gripper margins with translation or rotation margins.

\section*{4. \code{positive\_clip}}

\Diff{Difference: no centering and no z-score. Preserve raw token-level margins, but
only reinforce actions that the teacher scores above the old student.}

For each token:
\[
  \Adv_{t,b,k}
  =
  \min\left(c,\max\left(0,\Margin_{t,b,k}\right)\right).
\]

Interpretation:
\[
  \Margin_{t,b,k} > 0
  \Rightarrow
  \text{reinforce sampled action token},
\]
\[
  \Margin_{t,b,k} \le 0
  \Rightarrow
  \text{do not actively push against it}.
\]

This is conservative sampled imitation and may be useful when the teacher is strong on
the current task but unreliable for retained old-task behavior.

\section*{5. \code{teacher\_prob}}

\Diff{Difference: ignore the student margin and weight sampled actions by teacher
confidence. This is not full-distribution KL; it is sampled-action distillation.}

For each sampled action token:
\[
  w_{t,b,k}
  =
  \exp\left(
    \mathrm{clip}\left(
      \log \pi_{\mathrm{teacher}}(a_{t,b,k}\mid s_{t,b}),
      -20,
      0
    \right)
  \right).
\]

Normalize the scale over valid tokens:
\[
  \bar{w}
  =
  \frac{\sum_{t,b,k}\Mask_{t,b,k}w_{t,b,k}}
       {\max\left(1,\sum_{t,b,k}\Mask_{t,b,k}\right)}.
\]

Then:
\[
  \Adv_{t,b,k}
  =
  \frac{w_{t,b,k}}{\bar{w}+\epsilon}.
\]

Interpretation:
\[
  \text{high teacher probability}
  \Rightarrow
  \text{larger imitation weight};
  \qquad
  \text{low teacher probability}
  \Rightarrow
  \text{small update}.
\]

\section*{Continual-Learning Reading}

With a current-task SFT teacher, \code{group\_zscore}, \code{token\_zscore}, and
\code{action\_dim\_zscore} are still relative/centered objectives: some sampled
actions receive negative advantages even if the teacher is generally good. In contrast,
\code{positive\_clip} and \code{teacher\_prob} are more conservative imitation signals:
they mostly reinforce what the current-task teacher likes, but they do not by themselves
provide old-task retention signal. Forgetting prevention still requires replay,
old-task teacher/base regularization, EWC, or another source of old-task pressure.

\end{document}
